% !TeX program = lualatex
\documentclass[12pt,aspectratio=169]{beamer}
\usepackage[utf8]{inputenc}
\usepackage[T1]{fontenc}
\usepackage{amsmath,amsfonts,amssymb}
\usepackage{graphicx}
\usepackage{xcolor}
\usepackage{hyperref}
\usepackage{anyfontsize}
\usepackage{lmodern}
\usepackage{longtable}
\usepackage{array}
\usepackage{booktabs}

% ====== EXTREME FONT REDUCTION ======
\renewcommand{\tiny}{\fontsize{3.5}{4.5}\selectfont}
\renewcommand{\scriptsize}{\fontsize{4.5}{5.5}\selectfont}
\renewcommand{\footnotesize}{\fontsize{5.5}{6.5}\selectfont}
\renewcommand{\small}{\fontsize{6.5}{7.5}\selectfont}
\renewcommand{\normalsize}{\fontsize{7.5}{8.5}\selectfont}
\renewcommand{\large}{\fontsize{8.5}{9.5}\selectfont}
\renewcommand{\Large}{\fontsize{9.5}{10.5}\selectfont}
\renewcommand{\LARGE}{\fontsize{10.5}{12}\selectfont}
\renewcommand{\huge}{\fontsize{11.5}{13.5}\selectfont}
\renewcommand{\Huge}{\fontsize{13}{16}\selectfont}

% ====== COLORS ======
\definecolor{redcorrect}{RGB}{200,30,30}
\definecolor{bluecorrect}{RGB}{30,80,200}
\definecolor{greencheck}{RGB}{0,150,50}
\definecolor{lightgray}{RGB}{245,245,245}
\definecolor{darkgray}{RGB}{40,40,40}
\definecolor{softyellow}{RGB}{255,248,200}
\definecolor{deepblue}{RGB}{20,60,150}
\definecolor{darkred}{RGB}{150,20,20}
\definecolor{orangewarning}{RGB}{255,140,0}
\definecolor{correctgreen}{RGB}{0,120,40}
\definecolor{trapcolor}{RGB}{180,80,0}
\definecolor{purpleconcept}{RGB}{120,50,180}
\definecolor{tealhard}{RGB}{0,120,120}

\setbeamercolor{title}{fg=white,bg=redcorrect}
\setbeamercolor{frametitle}{fg=white,bg=redcorrect}
\setbeamercolor{block title}{fg=white,bg=bluecorrect}
\setbeamercolor{block body}{fg=darkgray,bg=lightgray}
\setbeamercolor{normal text}{fg=darkgray}
\usecolortheme{default}

% ====== CUSTOM COMMANDS ======
\newcommand{\correct}[1]{\textcolor{greencheck}{\textbf{\checkmark}} #1}
\newcommand{\keypoint}[1]{\textcolor{deepblue}{\textbf{#1}}}
\newcommand{\hardconcept}[1]{\textcolor{darkred}{\textbf{#1}}}
\newcommand{\examtraps}[1]{\textcolor{trapcolor}{\textbf{⚠ TRAP:}} #1}
\newcommand{\recallbox}[1]{\fcolorbox{deepblue}{blue!5}{\parbox{0.88\textwidth}{\centering \textbf{REMEMBER:} #1}}}
\newcommand{\stepbox}[1]{%
	\begin{center}
		\fcolorbox{darkgray}{white}{%
			\parbox{0.92\textwidth}{\vspace{0.03cm} #1 \vspace{0.03cm}}%
		}
	\end{center}
}
\newcommand{\answerbox}[1]{\fcolorbox{correctgreen}{green!8}{\parbox{0.88\textwidth}{\centering \textcolor{correctgreen}{\textbf{\checkmark\ CORRECT:}} #1}}}
\newcommand{\wronganswer}[1]{\textcolor{redcorrect}{\textbf{✘}} #1}
\newcommand{\trapbox}[1]{\fcolorbox{trapcolor}{orange!8}{\parbox{0.88\textwidth}{\centering \textcolor{trapcolor}{\textbf{⚠ EXAM TRAP:}} #1}}}
\newcommand{\keyrule}[1]{\textcolor{tealhard}{\textbf{🔑}} #1}

\begin{document}
	
	% ================================================================
	% SLIDE 1: TITLE
	% ================================================================
	\begin{frame}
		\centering
		\vspace{0.3cm}
		{\fontsize{18}{22}\selectfont \textbf{CAMS DOMAIN D}}\\[0.1cm]
		{\fontsize{14}{17}\selectfont \textbf{TOOLS \& TECHNOLOGIES TO FIGHT FINANCIAL CRIME}}\\[0.15cm]
		{\fontsize{9}{11}\selectfont \textit{ALL HARD QUESTIONS — COMPLETE TUTORIAL}}\\[0.08cm]
		\rule{4cm}{0.5pt}\\[0.08cm]
		{\fontsize{6}{7}\selectfont AI/ML Tools · RPA · Digital Onboarding · Data Validation · Data Lineage}\\[0.04cm]
		{\fontsize{6}{7}\selectfont Blockchain Analytics · Entity Resolution · Clustering · Batch Screening}\\[0.04cm]
		{\fontsize{6}{7}\selectfont Device Intelligence · Biometrics · Liveness Detection · Perpetual KYC}\\[0.04cm]
		{\fontsize{6}{7}\selectfont OSINT · Network Analysis · Cryptocurrency · Deepfake Detection}
		\vspace{0.2cm}
	\end{frame}
	
	% ================================================================
	% SLIDE 2: SANCTIONS LIST UPDATES
	% ================================================================
	\begin{frame}
		\frametitle{\fontsize{11}{13}\selectfont \textbf{Sanctions List Updates — The Aliases Trap}}
		\begin{block}{\fontsize{7}{9}\selectfont \textbf{The Hard Question — CAMS Exam Favorite}}
			\scriptsize
			In the context of list management, why is regular updating of sanctions lists critical to effective compliance operations?
		\end{block}
		\vspace{0.02cm}
		\answerbox{\large \textbf{Because sanctioned individuals frequently change affiliations and aliases.}}
		\vspace{0.02cm}
		\begin{block}{\fontsize{7}{9}\selectfont \textbf{Natural Explanation — Why You Got This Wrong}}
			\scriptsize
			\begin{itemize}
				\item \wronganswer{Because lists are only valid for 30 days} — \textbf{NO.} Lists are continuously updated.
				\item \wronganswer{Because new sanctions are added monthly} — \textbf{NO.} Updates are frequent and ad-hoc.
				\item \textbf{Sanctioned individuals:} Change \textbf{aliases}, \textbf{affiliations}, \textbf{locations}, and \textbf{connected entities}.
				\item \textbf{Regulators:} Update lists frequently — OFAC updates daily.
				\item \textbf{Risk:} Outdated lists = missed matches = regulatory violations.
				\item \textbf{Key principle:} Sanctions lists are \textbf{dynamic}, not static.
				\item \textbf{Exam Trap:} "30 days" or "monthly" — lists are updated \textbf{frequently and unpredictably}.
			\end{itemize}
		\end{block}
		\vspace{0.02cm}
		\recallbox{\scriptsize \textbf{Key Rule:} Sanctions lists update because \textbf{individuals change aliases/affiliations}.}
	\end{frame}
	
	% ================================================================
	% SLIDE 3: TRANSACTION MONITORING IMPROVEMENT
	% ================================================================
	\begin{frame}
		\frametitle{\fontsize{11}{13}\selectfont \textbf{Transaction Monitoring — The Behavioral Trap}}
		\begin{block}{\fontsize{7}{9}\selectfont \textbf{The Hard Question — CAMS Exam Favorite}}
			\scriptsize
			A bank's transaction monitoring system uses basic monetary thresholds to trigger alerts. Compliance personnel report high numbers of false positives and few true cases of financial crime. Which strategy would best improve the system?
		\end{block}
		\vspace{0.02cm}
		\answerbox{\large \textbf{Incorporate behavior and peer-group comparisons into the rules.}}
		\vspace{0.02cm}
		\begin{block}{\fontsize{7}{9}\selectfont \textbf{Natural Explanation — Why You Got This Wrong}}
			\scriptsize
			\begin{itemize}
				\item \wronganswer{Increase the number of rules to screen more transactions} — \textbf{NO.} More rules = more false positives.
				\item \textbf{Threshold-based systems:} Generate too many false positives.
				\item \textbf{Behavioral analysis:} Compares to \textbf{peer groups} — what's normal for that customer type.
				\item \textbf{Benefits:} Reduces false positives, detects \textbf{anomalies} rather than just high-value transactions.
				\item \textbf{Key principle:} Effective monitoring = \textbf{behavior + peer comparison}, not just thresholds.
				\item \textbf{Exam Trap:} "Increase rules" seems logical but creates \textbf{more noise}.
			\end{itemize}
		\end{block}
		\vspace{0.02cm}
		\recallbox{\scriptsize \textbf{Key Rule:} Improve monitoring = \textbf{behavioral + peer-group comparisons}.}
	\end{frame}
	
	% ================================================================
	% SLIDE 4: DYNAMIC DATA EXAMPLES
	% ================================================================
	\begin{frame}
		\frametitle{\fontsize{11}{13}\selectfont \textbf{Dynamic Data — The Examples Trap}}
		\begin{block}{\fontsize{7}{9}\selectfont \textbf{The Hard Question — CAMS Exam Favorite}}
			\scriptsize
			Which of the following are examples of dynamic data used in compliance controls? (Select Two.)
		\end{block}
		\vspace{0.02cm}
		\answerbox{\large \textbf{1. Real-time transaction geolocation. 2. Live adverse media feeds.}}
		\vspace{0.02cm}
		\begin{block}{\fontsize{7}{9}\selectfont \textbf{Natural Explanation — Why You Got This Wrong}}
			\scriptsize
			\begin{itemize}
				\item \wronganswer{Static customer KYC records} — \textbf{NO.} That's static, not dynamic.
				\item \textbf{Dynamic data:} Changes in real-time or near-real-time.
				\item \textbf{Geolocation:} Captures transaction location in real-time.
				\item \textbf{Adverse media feeds:} Continuously updated from news sources.
				\item \textbf{Other examples:} Transaction velocity, IP address, device fingerprint.
				\item \textbf{Key principle:} Dynamic = \textbf{real-time/continuous}. Static = \textbf{snapshot}.
				\item \textbf{Exam Trap:} "KYC records" are static — they are updated periodically, not dynamic.
			\end{itemize}
		\end{block}
		\vspace{0.02cm}
		\recallbox{\scriptsize \textbf{Key Rule:} Dynamic data = \textbf{real-time} (geolocation, live media feeds).}
	\end{frame}
	
	% ================================================================
	% SLIDE 5: GOVERNANCE FOR NEW TECHNOLOGY
	% ================================================================
	\begin{frame}
		\frametitle{\fontsize{11}{13}\selectfont \textbf{Governance — The Regulatory Misalignment Trap}}
		\begin{block}{\fontsize{7}{9}\selectfont \textbf{The Hard Question — CAMS Exam Favorite}}
			\scriptsize
			What risk may arise if governance is not established for newly implemented compliance technologies?
		\end{block}
		\vspace{0.02cm}
		\answerbox{\large \textbf{Outputs may misalign with regulatory expectations, diminishing overall effectiveness.}}
		\vspace{0.02cm}
		\begin{block}{\fontsize{7}{9}\selectfont \textbf{Natural Explanation — Why You Got This Wrong}}
			\scriptsize
			\begin{itemize}
				\item \wronganswer{The technology will stop working after 6 months} — \textbf{NO.} Technology doesn't "stop working" without governance.
				\item \wronganswer{Employees will refuse to use it} — \textbf{NO.} Governance is about compliance, not acceptance.
				\item \textbf{Without governance:} Systems may not align with regulatory requirements.
				\item \textbf{Governance:} Oversight, validation, testing, and documentation.
				\item \textbf{Key principle:} Governance ensures \textbf{alignment with regulatory expectations}.
				\item \textbf{Exam Trap:} "Technology stops working" is a trap — governance is about \textbf{compliance}, not functionality.
			\end{itemize}
		\end{block}
		\vspace{0.02cm}
		\recallbox{\scriptsize \textbf{Key Rule:} Governance = \textbf{alignment with regulatory expectations}.}
	\end{frame}
	
	% ================================================================
	% SLIDE 6: EXTERNAL DATA FOR ONGOING MONITORING
	% ================================================================
	\begin{frame}
		\frametitle{\fontsize{11}{13}\selectfont \textbf{External Data — The Structured Trap}}
		\begin{block}{\fontsize{7}{9}\selectfont \textbf{The Hard Question — CAMS Exam Favorite}}
			\scriptsize
			Which actions would help a compliance team improve the usefulness of external data for ongoing monitoring? (Select Two.)
		\end{block}
		\vspace{0.02cm}
		\answerbox{\large \textbf{1. Create a dedicated data pipeline to continually scan public court and business registry records. 2. Standardize incoming external data from various formats into a structured schema.}}
		\vspace{0.02cm}
		\begin{block}{\fontsize{7}{9}\selectfont \textbf{Natural Explanation — Why You Got This Wrong}}
			\scriptsize
			\begin{itemize}
				\item \wronganswer{Remove structured internal data to prevent conflicts with external lists} — \textbf{NO.} You need \textbf{both} internal and external data.
				\item \textbf{Data pipeline:} Continuously ingest external data.
				\item \textbf{Standardization:} Convert various formats (PDF, JSON, XML) into a \textbf{structured schema}.
				\item \textbf{Key principle:} External data usefulness = \textbf{continuous ingestion + standardization}.
				\item \textbf{Exam Trap:} "Remove internal data" is dangerous — you need integration, not removal.
			\end{itemize}
		\end{block}
		\vspace{0.02cm}
		\recallbox{\scriptsize \textbf{Key Rule:} External data = \textbf{continuous pipeline + standardization}.}
	\end{frame}
	
	% ================================================================
	% SLIDE 7: CUSTOMER RISK SCORING ENGINE
	% ================================================================
	\begin{frame}
		\frametitle{\fontsize{11}{13}\selectfont \textbf{Risk Scoring — The Success Interpretation Trap}}
		\begin{block}{\fontsize{7}{9}\selectfont \textbf{The Hard Question — CAMS Exam Favorite}}
			\scriptsize
			An organization is piloting a customer-risk scoring engine that combines internal and external data. After the first quarter, high-risk clients that were previously undetected are flagged. How should this result be interpreted?
		\end{block}
		\vspace{0.02cm}
		\answerbox{\large \textbf{This may indicate successful performance by identifying missed risk exposure.}}
		\vspace{0.02cm}
		\begin{block}{\fontsize{7}{9}\selectfont \textbf{Natural Explanation — Why You Got This Wrong}}
			\scriptsize
			\begin{itemize}
				\item \wronganswer{The engine is over-sensitive and overrates risk} — \textbf{NO.} New flags don't mean false positives.
				\item \textbf{Risk scoring engine:} Combines internal and external data.
				\item \textbf{Previously undetected:} May have been missed by previous system.
				\item \textbf{Key principle:} New flags = \textbf{may indicate success} at identifying previously missed risk.
				\item \textbf{Key principle:} Additional flags = \textbf{more effective detection}, not necessarily over-sensitivity.
				\item \textbf{Exam Trap:} "Over-sensitive" assumes false positives — but new flags may be \textbf{true positives}.
			\end{itemize}
		\end{block}
		\vspace{0.02cm}
		\recallbox{\scriptsize \textbf{Key Rule:} New flags = \textbf{possible success at identifying missed risk}.}
	\end{frame}
	
	% ================================================================
	% SLIDE 8: KYC ONBOARDING IMPROVEMENTS
	% ================================================================
	\begin{frame}
		\frametitle{\fontsize{11}{13}\selectfont \textbf{KYC Onboarding — The Automation Trap}}
		\begin{block}{\fontsize{7}{9}\selectfont \textbf{The Hard Question — CAMS Exam Favorite}}
			\scriptsize
			An international bank is evaluating KYC process improvements. Its current onboarding process suffers from high error rates in data entry and document verification. Based on recent technological advancements, which is most likely to be improved?
		\end{block}
		\vspace{0.02cm}
		\answerbox{\large \textbf{Automated data entry and document verification processes.}}
		\vspace{0.02cm}
		\begin{block}{\fontsize{7}{9}\selectfont \textbf{Natural Explanation — Why You Got This Wrong}}
			\scriptsize
			\begin{itemize}
				\item \wronganswer{Manual verification of all documents} — \textbf{NO.} Manual processes don't reduce errors.
				\item \wronganswer{Eliminating all data entry} — \textbf{NO.} You still need data, just automated.
				\item \textbf{Automated data entry:} OCR, AI-powered document extraction.
				\item \textbf{Document verification:} Biometric checks, ID validation against government databases.
				\item \textbf{Key principle:} Technology reduces \textbf{manual errors} through \textbf{automation}.
				\item \textbf{Exam Trap:} "Manual verification" is the opposite of improvement — automation is the solution.
			\end{itemize}
		\end{block}
		\vspace{0.02cm}
		\recallbox{\scriptsize \textbf{Key Rule:} KYC improvement = \textbf{automated data entry + document verification}.}
	\end{frame}
	
	% ================================================================
	% SLIDE 9: DOCUMENTATION FOR OSINT AND AI
	% ================================================================
	\begin{frame}
		\frametitle{\fontsize{11}{13}\selectfont \textbf{OSINT/AI Documentation — The Feature Selection Trap}}
		\begin{block}{\fontsize{7}{9}\selectfont \textbf{The Hard Question — CAMS Exam Favorite}}
			\scriptsize
			Which describe components of proper documentation practices for OSINT and AI-driven monitoring? (Select Two.)
		\end{block}
		\vspace{0.02cm}
		\answerbox{\large \textbf{1. Feature selection rationale and training dataset details. 2. Explanations for model output assumptions using technical metrics.}}
		\vspace{0.02cm}
		\begin{block}{\fontsize{7}{9}\selectfont \textbf{Natural Explanation — Why You Got This Wrong}}
			\scriptsize
			\begin{itemize}
				\item \wronganswer{Audit logs of all internal system errors} — \textbf{NO.} Error logs are operational, not documentation for AI/OSINT.
				\item \textbf{Feature selection rationale:} Why each feature was chosen for the model.
				\item \textbf{Training dataset details:} What data was used to train the model.
				\item \textbf{Model output assumptions:} How the model makes decisions.
				\item \textbf{Key principle:} AI/OSINT documentation = \textbf{explainability + transparency}.
				\item \textbf{Exam Trap:} "Audit logs" are operational — AI documentation is about \textbf{model decisions}.
			\end{itemize}
		\end{block}
		\vspace{0.02cm}
		\recallbox{\scriptsize \textbf{Key Rule:} AI documentation = \textbf{feature selection + training data + output explanations}.}
	\end{frame}
	
	% ================================================================
	% SLIDE 10: CLUSTERING ANALYSIS — MEANINGFUL CONNECTIONS
	% ================================================================
	\begin{frame}
		\frametitle{\fontsize{11}{13}\selectfont \textbf{Clustering — The Connection Trap}}
		\begin{block}{\fontsize{7}{9}\selectfont \textbf{The Hard Question — CAMS Exam Favorite}}
			\scriptsize
			During a clustering analysis, which combination would most suggest a meaningful connection between accounts?
		\end{block}
		\vspace{0.02cm}
		\answerbox{\large \textbf{Same employer address, similar transaction values, interlinked payments.}}
		\vspace{0.02cm}
		\begin{block}{\fontsize{7}{9}\selectfont \textbf{Natural Explanation — Why You Got This Wrong}}
			\scriptsize
			\begin{itemize}
				\item \wronganswer{Same bank, different branches} — \textbf{NO.} That's not a connection — just same bank.
				\item \wronganswer{Similar account balance ranges} — \textbf{NO.} That's just coincidental.
				\item \textbf{Same employer address:} Shared physical location.
				\item \textbf{Similar transaction values:} Pattern of similar amounts.
				\item \textbf{Interlinked payments:} Payments between accounts.
				\item \textbf{Key principle:} Meaningful connection = \textbf{multiple shared attributes + transaction links}.
				\item \textbf{Exam Trap:} "Same bank" sounds like a connection but is too broad.
			\end{itemize}
		\end{block}
		\vspace{0.02cm}
		\recallbox{\scriptsize \textbf{Key Rule:} Meaningful connection = \textbf{shared attributes + transaction links}.}
	\end{frame}
	
	% ================================================================
	% SLIDE 11: HUMAN-IN-THE-LOOP TESTING
	% ================================================================
	\begin{frame}
		\frametitle{\fontsize{11}{13}\selectfont \textbf{Human-in-the-Loop — The Explainability Trap}}
		\begin{block}{\fontsize{7}{9}\selectfont \textbf{The Hard Question — CAMS Exam Favorite}}
			\scriptsize
			Which best describe human-in-the-loop testing in AI-based AML systems? (Select Two.)
		\end{block}
		\vspace{0.02cm}
		\answerbox{\large \textbf{1. Experts assess model decisions for accuracy. 2. It ensures explainability of model outputs.}}
		\vspace{0.02cm}
		\begin{block}{\fontsize{7}{9}\selectfont \textbf{Natural Explanation — Why You Got This Wrong}}
			\scriptsize
			\begin{itemize}
				\item \wronganswer{It eliminates all false positives} — \textbf{NO.} Human review doesn't eliminate false positives.
				\item \wronganswer{It replaces the AI model} — \textbf{NO.} It's an oversight mechanism, not a replacement.
				\item \textbf{Experts assess model decisions:} Validate accuracy of AI outputs.
				\item \textbf{Ensures explainability:} Humans can explain why a decision was made.
				\item \textbf{Key principle:} Human-in-the-loop = \textbf{oversight + explainability}.
				\item \textbf{Exam Trap:} "Eliminates false positives" is unrealistic — humans still make errors.
			\end{itemize}
		\end{block}
		\vspace{0.02cm}
		\recallbox{\scriptsize \textbf{Key Rule:} Human-in-the-loop = \textbf{expert assessment + explainability}.}
	\end{frame}
	
	% ================================================================
	% SLIDE 12: ENTITY RESOLUTION CHALLENGES
	% ================================================================
	\begin{frame}
		\frametitle{\fontsize{11}{13}\selectfont \textbf{Entity Resolution — The Connection Trap}}
		\begin{block}{\fontsize{7}{9}\selectfont \textbf{The Hard Question — CAMS Exam Favorite}}
			\scriptsize
			A large global financial institution is implementing an entity resolution system to improve its compliance processes. What compliance challenge might be revealed after implementation?
		\end{block}
		\vspace{0.02cm}
		\answerbox{\large \textbf{Previously undetected connections between seemingly independent businesses with common ownership.}}
		\vspace{0.02cm}
		\begin{block}{\fontsize{7}{9}\selectfont \textbf{Natural Explanation — Why You Got This Wrong}}
			\scriptsize
			\begin{itemize}
				\item \wronganswer{Incompatibility between entity resolution algorithms and existing systems} — This is an \textbf{implementation challenge}, not a compliance revelation.
				\item \textbf{Entity resolution:} Identifies relationships between entities.
				\item \textbf{Reveals:} Hidden connections — common ownership, shared addresses, related parties.
				\item \textbf{Key principle:} Entity resolution = \textbf{reveals hidden connections}.
				\item \textbf{Exam Trap:} "Incompatibility" is a technical issue — the question asks about compliance challenges \textbf{revealed} after implementation.
			\end{itemize}
		\end{block}
		\vspace{0.02cm}
		\recallbox{\scriptsize \textbf{Key Rule:} Entity resolution reveals \textbf{previously undetected connections}.}
	\end{frame}
	
	% ================================================================
	% SLIDE 13: BEHAVIORAL MONITORING GOVERNANCE
	% ================================================================
	\begin{frame}
		\frametitle{\fontsize{11}{13}\selectfont \textbf{Behavioral Monitoring — The Governance Trap}}
		\begin{block}{\fontsize{7}{9}\selectfont \textbf{The Hard Question — CAMS Exam Favorite}}
			\scriptsize
			Which should inform governance committees as part of oversight of behavioral monitoring technologies? (Select Two.)
		\end{block}
		\vspace{0.02cm}
		\answerbox{\large \textbf{1. Alert thresholds compared to actual suspicious activity reports. 2. Results of periodic tuning reviews.}}
		\vspace{0.02cm}
		\begin{block}{\fontsize{7}{9}\selectfont \textbf{Natural Explanation — Why You Got This Wrong}}
			\scriptsize
			\begin{itemize}
				\item \wronganswer{Records of minimum training compliance} — \textbf{NO.} Training is operational, not governance oversight.
				\item \wronganswer{Metrics monitoring speed of alert resolution} — \textbf{NO.} Speed is operational, not governance.
				\item \textbf{Alert thresholds vs. SARs:} Validates if thresholds are effective.
				\item \textbf{Periodic tuning reviews:} Ensures system remains calibrated.
				\item \textbf{Key principle:} Governance = \textbf{effectiveness + calibration}, not operational metrics.
				\item \textbf{Exam Trap:} "Training compliance" and "alert speed" are operational — governance is about \textbf{effectiveness}.
			\end{itemize}
		\end{block}
		\vspace{0.02cm}
		\recallbox{\scriptsize \textbf{Key Rule:} Behavioral monitoring governance = \textbf{thresholds vs. SARs + tuning reviews}.}
	\end{frame}
	
	% ================================================================
	% SLIDE 14: OSINT QUALITY GOVERNANCE
	% ================================================================
	\begin{frame}
		\frametitle{\fontsize{11}{13}\selectfont \textbf{OSINT Quality — The Source Assessment Trap}}
		\begin{block}{\fontsize{7}{9}\selectfont \textbf{The Hard Question — CAMS Exam Favorite}}
			\scriptsize
			An institution uses OSINT to enrich profiles of high-risk customers. Recently, repeated false positives were traced to poor-quality web scraping. What should be the institution's next governance priority?
		\end{block}
		\vspace{0.02cm}
		\answerbox{\large \textbf{Reassess OSINT input sources and document reliability assessments.}}
		\vspace{0.02cm}
		\begin{block}{\fontsize{7}{9}\selectfont \textbf{Natural Explanation — Why You Got This Wrong}}
			\scriptsize
			\begin{itemize}
				\item \wronganswer{Increase the volume of web scraping} — \textbf{NO.} More scraping = more poor-quality data.
				\item \wronganswer{Stop using OSINT entirely} — \textbf{NO.} OSINT is valuable; need to improve quality.
				\item \textbf{Reassess input sources:} Evaluate reliability of each data source.
				\item \textbf{Document reliability assessments:} Record quality of each source.
				\item \textbf{Key principle:} OSINT quality = \textbf{source reliability assessment}.
				\item \textbf{Exam Trap:} "Increase scraping" assumes quantity = quality — \textbf{incorrect.}
			\end{itemize}
		\end{block}
		\vspace{0.02cm}
		\recallbox{\scriptsize \textbf{Key Rule:} OSINT quality = \textbf{reassess sources + document reliability}.}
	\end{frame}
	
	% ================================================================
	% SLIDE 15: RISK-BASED AFC TECHNOLOGY SELECTION
	% ================================================================
	\begin{frame}
		\frametitle{\fontsize{11}{13}\selectfont \textbf{Risk-Based Tech Selection — The Approach Trap}}
		\begin{block}{\fontsize{7}{9}\selectfont \textbf{The Hard Question — CAMS Exam Favorite}}
			\scriptsize
			Which actions support a risk-based approach in selecting AFC technologies? (Select Three.)
		\end{block}
		\vspace{0.02cm}
		\answerbox{\large \textbf{1. Assess the technology's ability to detect typologies relevant to the institution's risk profile. 2. Evaluate the vendor's history and reputation. 3. Align technology capabilities with the institution's compliance strategy.}}
		\vspace{0.02cm}
		\begin{block}{\fontsize{7}{9}\selectfont \textbf{Natural Explanation — Why You Got This Wrong}}
			\scriptsize
			\begin{itemize}
				\item \wronganswer{Select the cheapest technology available} — \textbf{NO.} Cost is not the primary risk-based factor.
				\item \wronganswer{Choose the most advanced technology regardless of risk} — \textbf{NO.} Must align with risk profile.
				\item \textbf{Risk-based approach:} Match technology to \textbf{institution's risk profile}.
				\item \textbf{Assess typology detection:} Can it detect relevant financial crime types?
				\item \textbf{Vendor evaluation:} Reputation, history, reliability.
				\item \textbf{Align with strategy:} Technology must fit compliance goals.
				\item \textbf{Key principle:} Risk-based = \textbf{fit for purpose}, not cheapest or most advanced.
				\item \textbf{Exam Trap:} "Cheapest" and "most advanced" ignore risk alignment.
			\end{itemize}
		\end{block}
		\vspace{0.02cm}
		\recallbox{\scriptsize \textbf{Key Rule:} Risk-based tech selection = \textbf{fit for risk profile + vendor evaluation + strategy alignment}.}
	\end{frame}
	
	% ================================================================
	% SLIDE 16: MACHINE LEARNING BENEFITS
	% ================================================================
	\begin{frame}
		\frametitle{\fontsize{11}{13}\selectfont \textbf{Machine Learning — The Improvement Trap}}
		\begin{block}{\fontsize{7}{9}\selectfont \textbf{The Hard Question — CAMS Exam Favorite}}
			\scriptsize
			Which are benefits of using machine learning for customer screening? (Select Two.)
		\end{block}
		\vspace{0.02cm}
		\answerbox{\large \textbf{1. Continuous learning from recent screening results. 2. Improved detection accuracy over time.}}
		\vspace{0.02cm}
		\begin{block}{\fontsize{7}{9}\selectfont \textbf{Natural Explanation — Why You Got This Wrong}}
			\scriptsize
			\begin{itemize}
				\item \wronganswer{Reduced compliance staff headcount} — \textbf{NO.} ML augments staff, doesn't reduce headcount.
				\item \textbf{Continuous learning:} ML models improve as they process more data.
				\item \textbf{Improved accuracy:} Detection accuracy improves over time with training.
				\item \textbf{Key principle:} ML = \textbf{continuous learning + improved accuracy}.
				\item \textbf{Exam Trap:} "Reduced headcount" is a common misconception — ML \textbf{augments}, not replaces.
			\end{itemize}
		\end{block}
		\vspace{0.02cm}
		\recallbox{\scriptsize \textbf{Key Rule:} ML benefits = \textbf{continuous learning + improved accuracy}.}
	\end{frame}
	
	% ================================================================
	% SLIDE 17: CARD-BASED PRODUCT RISKS
	% ================================================================
	\begin{frame}
		\frametitle{\fontsize{11}{13}\selectfont \textbf{Card Products — The Anonymity Trap}}
		\begin{block}{\fontsize{7}{9}\selectfont \textbf{The Hard Question — CAMS Exam Favorite}}
			\scriptsize
			Which represent common risk factors associated with card-based products? (Select Three.)
		\end{block}
		\vspace{0.02cm}
		\answerbox{\large \textbf{1. Anonymity through prepaid cards. 2. High volumes of rapid transactions. 3. Use by customers in high-risk sectors.}}
		\vspace{0.02cm}
		\begin{block}{\fontsize{7}{9}\selectfont \textbf{Natural Explanation — Why You Got This Wrong}}
			\scriptsize
			\begin{itemize}
				\item \wronganswer{Cards are always subject to low fraud risk} — \textbf{NO.} Cards have significant fraud risk.
				\item \textbf{Prepaid card anonymity:} Can be loaded without full KYC.
				\item \textbf{High volume rapid transactions:} Structured transactions to avoid reporting.
				\item \textbf{High-risk sectors:} Cards used by customers in high-risk industries.
				\item \textbf{Key principle:} Card risks = \textbf{anonymity + velocity + customer type}.
				\item \textbf{Exam Trap:} "Low fraud risk" is a trap — cards are high-risk products.
			\end{itemize}
		\end{block}
		\vspace{0.02cm}
		\recallbox{\scriptsize \textbf{Key Rule:} Card risks = \textbf{anonymity + rapid transactions + high-risk customers}.}
	\end{frame}
	
	% ================================================================
	% SLIDE 18: DATA COVERAGE GAP ASSESSMENT
	% ================================================================
	\begin{frame}
		\frametitle{\fontsize{11}{13}\selectfont \textbf{Data Gap Assessment — The Missing Identifiers Trap}}
		\begin{block}{\fontsize{7}{9}\selectfont \textbf{The Hard Question — CAMS Exam Favorite}}
			\scriptsize
			Which challenges can be identified through a data coverage and gap assessment? (Select Two.)
		\end{block}
		\vspace{0.02cm}
		\answerbox{\large \textbf{1. Missing identifiers in customer profiles. 2. Inconsistent use of country codes across systems.}}
		\vspace{0.02cm}
		\begin{block}{\fontsize{7}{9}\selectfont \textbf{Natural Explanation — Why You Got This Wrong}}
			\scriptsize
			\begin{itemize}
				\item \wronganswer{Misalignment between transaction monitoring scenarios and typologies} — \textbf{NO.} That's a \textbf{scenario design issue}, not a data gap.
				\item \textbf{Missing identifiers:} Incomplete KYC data (missing DOB, ID numbers, addresses).
				\item \textbf{Inconsistent country codes:} Different systems use different codes (US vs USA vs 840).
				\item \textbf{Key principle:} Data gap assessment = \textbf{data completeness + consistency}.
				\item \textbf{Exam Trap:} "Misalignment with typologies" is a \textbf{scenario issue}, not a data gap.
			\end{itemize}
		\end{block}
		\vspace{0.02cm}
		\recallbox{\scriptsize \textbf{Key Rule:} Data gaps = \textbf{missing identifiers + inconsistent formats}.}
	\end{frame}
	
	% ================================================================
	% SLIDE 19: DNFBP SUPERVISION
	% ================================================================
	\begin{frame}
		\frametitle{\fontsize{11}{13}\selectfont \textbf{DNFBP — The Supervision Trap}}
		\begin{block}{\fontsize{7}{9}\selectfont \textbf{The Hard Question — CAMS Exam Favorite}}
			\scriptsize
			Following a mutual evaluation, a jurisdiction is asked to improve its supervision of designated non-financial businesses and professions (DNFBPs). Which measure would most likely satisfy this request?
		\end{block}
		\vspace{0.02cm}
		\answerbox{\large \textbf{Expand AML/CFT requirements to include lawyers, accountants, and real estate agents.}}
		\vspace{0.02cm}
		\begin{block}{\fontsize{7}{9}\selectfont \textbf{Natural Explanation — Why You Got This Wrong}}
			\scriptsize
			\begin{itemize}
				\item \wronganswer{Increase penalties for financial institutions} — \textbf{NO.} DNFBPs are non-financial, not financial institutions.
				\item \textbf{DNFBPs:} Lawyers, accountants, real estate agents, trust companies, casinos.
				\item \textbf{Supervision:} Must include AML/CFT requirements for these professions.
				\item \textbf{Key principle:} DNFBP supervision = \textbf{extend AML/CFT requirements} to non-financial sectors.
				\item \textbf{Exam Trap:} "Increase penalties for banks" misses the point — DNFBPs are the issue.
			\end{itemize}
		\end{block}
		\vspace{0.02cm}
		\recallbox{\scriptsize \textbf{Key Rule:} DNFBP supervision = \textbf{extend AML/CFT to lawyers, accountants, real estate}.}
	\end{frame}
	
	% ================================================================
	% SLIDE 20: PSP MULTI-JURISDICTION
	% ================================================================
	\begin{frame}
		\frametitle{\fontsize{11}{13}\selectfont \textbf{PSP Strategy — The Customization Trap}}
		\begin{block}{\fontsize{7}{9}\selectfont \textbf{The Hard Question — CAMS Exam Favorite}}
			\scriptsize
			A PSP is planning to offer new services targeting small merchants in multiple countries, some lacking mature AML frameworks. To maintain compliance and mitigate risk, which strategy is most appropriate?
		\end{block}
		\vspace{0.02cm}
		\answerbox{\large \textbf{Customize AML policy and onboarding procedures per jurisdiction and product.}}
		\vspace{0.02cm}
		\begin{block}{\fontsize{7}{9}\selectfont \textbf{Natural Explanation — Why You Got This Wrong}}
			\scriptsize
			\begin{itemize}
				\item \wronganswer{Apply home-country AML standards across all jurisdictions} — \textbf{NO.} Home standards may not meet local requirements.
				\item \textbf{Customization:} Each jurisdiction has different AML/CFT requirements.
				\item \textbf{Product-specific:} Different products have different risk profiles.
				\item \textbf{Key principle:} Multi-jurisdiction = \textbf{customize per jurisdiction and product}.
				\item \textbf{Exam Trap:} "Home-country standards" oversimplifies — must comply with \textbf{local} laws.
			\end{itemize}
		\end{block}
		\vspace{0.02cm}
		\recallbox{\scriptsize \textbf{Key Rule:} PSP strategy = \textbf{customize per jurisdiction and product}.}
	\end{frame}
	
	% ================================================================
	% SLIDE 21: PREPAID PRODUCT RESPONSE
	% ================================================================
	\begin{frame}
		\frametitle{\fontsize{11}{13}\selectfont \textbf{Prepaid Product — The Reassessment Trap}}
		\begin{block}{\fontsize{7}{9}\selectfont \textbf{The Hard Question — CAMS Exam Favorite}}
			\scriptsize
			An international financial institution launches a new prepaid product and later identifies a pattern of high-frequency transactions from high-risk jurisdictions shortly after launch. What should the institution do FIRST?
		\end{block}
		\vspace{0.02cm}
		\answerbox{\large \textbf{Reassess the product's associated risks and adjust transaction thresholds.}}
		\vspace{0.02cm}
		\begin{block}{\fontsize{7}{9}\selectfont \textbf{Natural Explanation — Why You Got This Wrong}}
			\scriptsize
			\begin{itemize}
				\item \wronganswer{Notify law enforcement and suspend the product} — \textbf{NO.} That's too drastic without assessment.
				\item \textbf{FIRST step:} Reassess the product's risk profile.
				\item \textbf{Adjust thresholds:} Lower transaction limits for high-risk jurisdictions.
				\item \textbf{Key principle:} First step = \textbf{risk reassessment}, not suspension.
				\item \textbf{Exam Trap:} "Notify law enforcement" is premature — must assess first.
			\end{itemize}
		\end{block}
		\vspace{0.02cm}
		\recallbox{\scriptsize \textbf{Key Rule:} Prepaid product issue = \textbf{reassess risks + adjust thresholds} FIRST.}
	\end{frame}
	
	% ================================================================
	% SLIDE 22: CLUSTERING CHARACTERISTICS
	% ================================================================
	\begin{frame}
		\frametitle{\fontsize{11}{13}\selectfont \textbf{Clustering Characteristics — The Connections Trap}}
		\begin{block}{\fontsize{7}{9}\selectfont \textbf{The Hard Question — CAMS Exam Favorite}}
			\scriptsize
			Which are characteristics of clustering in compliance analytics? (Select Two.)
		\end{block}
		\vspace{0.02cm}
		\answerbox{\large \textbf{1. It groups entities by shared behavioral characteristics. 2. It identifies connections across multiple transaction types.}}
		\vspace{0.02cm}
		\begin{block}{\fontsize{7}{9}\selectfont \textbf{Natural Explanation — Why You Got This Wrong}}
			\scriptsize
			\begin{itemize}
				\item \wronganswer{It excludes outliers to reduce financial crime risk} — \textbf{NO.} Outliers are often the \textbf{most suspicious}.
				\item \textbf{Clustering:} Groups entities with similar behaviors.
				\item \textbf{Connections across transaction types:} Links different transaction types (wire, ACH, card) that show similar patterns.
				\item \textbf{Key principle:} Clustering = \textbf{shared behavior + cross-transaction connections}.
				\item \textbf{Exam Trap:} "Excluding outliers" misses suspicious activity — outliers are key.
			\end{itemize}
		\end{block}
		\vspace{0.02cm}
		\recallbox{\scriptsize \textbf{Key Rule:} Clustering = \textbf{shared behaviors + cross-transaction connections}.}
	\end{frame}
	
	% ================================================================
	% SLIDE 23: DATA VISUALIZATION TOOL ISSUE
	% ================================================================
	\begin{frame}
		\frametitle{\fontsize{11}{13}\selectfont \textbf{Data Visualization — The Speed Trap}}
		\begin{block}{\fontsize{7}{9}\selectfont \textbf{The Hard Question — CAMS Exam Favorite}}
			\scriptsize
			During a quarterly performance audit of a data visualization-based analysis tool, investigators note that investigative resolution time decreased by 40\%, but false negatives have increased. What conclusion is most appropriate?
		\end{block}
		\vspace{0.02cm}
		\answerbox{\large \textbf{Analysts are resolving too quickly and overlooking weaker signals.}}
		\vspace{0.02cm}
		\begin{block}{\fontsize{7}{9}\selectfont \textbf{Natural Explanation — Why You Got This Wrong}}
			\scriptsize
			\begin{itemize}
				\item \wronganswer{The tool is too visual and fails to produce structured insights} — \textbf{NO.} Visualization tools are designed to produce insights.
				\item \textbf{Decreased resolution time:} Faster is not always better.
				\item \textbf{Increased false negatives:} Missing suspicious activity.
				\item \textbf{Key principle:} Speed without accuracy = \textbf{overlooking signals}.
				\item \textbf{Exam Trap:} "Too visual" is a red herring — the issue is \textbf{analysis speed vs. accuracy}.
			\end{itemize}
		\end{block}
		\vspace{0.02cm}
		\recallbox{\scriptsize \textbf{Key Rule:} Faster resolution + more false negatives = \textbf{overlooking weaker signals}.}
	\end{frame}
	
	% ================================================================
	% SLIDE 24: DATA VALIDATION OBJECTIVES
	% ================================================================
	\begin{frame}
		\frametitle{\fontsize{11}{13}\selectfont \textbf{Data Validation — The Objectives Trap}}
		\begin{block}{\fontsize{7}{9}\selectfont \textbf{The Hard Question — CAMS Exam Favorite}}
			\scriptsize
			Which are typical objectives of data validation and testing in AFC systems? (Select Three.)
		\end{block}
		\vspace{0.02cm}
		\answerbox{\large \textbf{1. Ensuring data inputs reflect actual customer and transaction behaviors. 2. Verifying automated controls are functioning correctly. 3. Testing whether controls generate alerts under appropriate risk conditions.}}
		\vspace{0.02cm}
		\begin{block}{\fontsize{7}{9}\selectfont \textbf{Natural Explanation — Why You Got This Wrong}}
			\scriptsize
			\begin{itemize}
				\item \wronganswer{Confirming reports are generated in a format preferred by third-party vendors} — \textbf{NO.} That's operational, not validation.
				\item \textbf{Data inputs reflect behavior:} Are we capturing the right data?
				\item \textbf{Controls functioning correctly:} Are the systems working as designed?
				\item \textbf{Alerts under risk conditions:} Do we get alerts when we should?
				\item \textbf{Key principle:} Validation = \textbf{accuracy + functionality + appropriate alerts}.
				\item \textbf{Exam Trap:} "Vendor format preference" is operational — not a validation objective.
			\end{itemize}
		\end{block}
		\vspace{0.02cm}
		\recallbox{\scriptsize \textbf{Key Rule:} Data validation = \textbf{accuracy + functionality + appropriate alerts}.}
	\end{frame}
	
	% ================================================================
	% SLIDE 25: DATA LINEAGE
	% ================================================================
	\begin{frame}
		\frametitle{\fontsize{11}{13}\selectfont \textbf{Data Lineage — The Transformation Trap}}
		\begin{block}{\fontsize{7}{9}\selectfont \textbf{The Hard Question — CAMS Exam Favorite}}
			\scriptsize
			Which support transparency and auditability through data lineage? (Select Two.)
		\end{block}
		\vspace{0.02cm}
		\answerbox{\large \textbf{1. Maintaining history of data transformations. 2. Capturing source of each data field.}}
		\vspace{0.02cm}
		\begin{block}{\fontsize{7}{9}\selectfont \textbf{Natural Explanation — Why You Got This Wrong}}
			\scriptsize
			\begin{itemize}
				\item \wronganswer{Documenting rule development timelines} — \textbf{NO.} That's rule documentation, not data lineage.
				\item \textbf{Data transformations:} Records how data was changed.
				\item \textbf{Source of data:} Where each field originated.
				\item \textbf{Key principle:} Data lineage = \textbf{transformations + source}.
				\item \textbf{Exam Trap:} "Rule development timelines" is \textbf{not} data lineage.
			\end{itemize}
		\end{block}
		\vspace{0.02cm}
		\recallbox{\scriptsize \textbf{Key Rule:} Data lineage = \textbf{transformations history + source of each field}.}
	\end{frame}
	
	% ================================================================
	% SLIDE 26: REAL-TIME TERRORIST FINANCING DETECTION
	% ================================================================
	\begin{frame}
		\frametitle{\fontsize{11}{13}\selectfont \textbf{Real-Time Detection — The Anomaly Trap}}
		\begin{block}{\fontsize{7}{9}\selectfont \textbf{The Hard Question — CAMS Exam Favorite}}
			\scriptsize
			A fintech startup is piloting a new payment screening program aimed at preventing terrorist financing through real-time detection. What technical feature should they prioritize to ensure early pattern recognition?
		\end{block}
		\vspace{0.02cm}
		\answerbox{\large \textbf{Machine learning systems trained on prior known suspicious behavior to detect real-time anomalies.}}
		\vspace{0.02cm}
		\begin{block}{\fontsize{7}{9}\selectfont \textbf{Natural Explanation — Why You Got This Wrong}}
			\scriptsize
			\begin{itemize}
				\item \wronganswer{Manual review by human analysts} — \textbf{NO.} Manual review cannot keep up with real-time.
				\item \wronganswer{A rule-based system with predefined thresholds} — \textbf{NO.} Rules don't detect new patterns.
				\item \textbf{ML trained on known suspicious behavior:} Learns patterns of terrorist financing.
				\item \textbf{Real-time anomalies:} Detects suspicious activity as it occurs.
				\item \textbf{Key principle:} Real-time detection = \textbf{ML trained on known suspicious behavior}.
				\item \textbf{Exam Trap:} "Rules" don't adapt — ML learns new patterns.
			\end{itemize}
		\end{block}
		\vspace{0.02cm}
		\recallbox{\scriptsize \textbf{Key Rule:} Real-time detection = \textbf{ML trained on known suspicious behavior}.}
	\end{frame}
	
	% ================================================================
	% SLIDE 27: DATA STANDARDIZATION
	% ================================================================
	\begin{frame}
		\frametitle{\fontsize{11}{13}\selectfont \textbf{Data Standardization — The Mapping Schema Trap}}
		\begin{block}{\fontsize{7}{9}\selectfont \textbf{The Hard Question — CAMS Exam Favorite}}
			\scriptsize
			A global institution operates multiple platforms with varying labels for "beneficiary" data, causing confusion during system transfers. What integrated solution should it prioritize?
		\end{block}
		\vspace{0.02cm}
		\answerbox{\large \textbf{Developing a standardized mapping schema linked to a data dictionary.}}
		\vspace{0.02cm}
		\begin{block}{\fontsize{7}{9}\selectfont \textbf{Natural Explanation — Why You Got This Wrong}}
			\scriptsize
			\begin{itemize}
				\item \wronganswer{Using the same system for all platforms} — \textbf{NO.} Different platforms have different data structures.
				\item \textbf{Data dictionary:} Defines what each field means.
				\item \textbf{Mapping schema:} Shows how data maps to the dictionary.
				\item \textbf{Key principle:} Standardization = \textbf{dictionary + mapping}.
				\item \textbf{Exam Trap:} "Same system" ignores data structure differences.
			\end{itemize}
		\end{block}
		\vspace{0.02cm}
		\recallbox{\scriptsize \textbf{Key Rule:} Data standardization = \textbf{data dictionary + mapping schema}.}
	\end{frame}
	
	% ================================================================
	% SLIDE 28: DIGITAL ONBOARDING BENEFITS
	% ================================================================
	\begin{frame}
		\frametitle{\fontsize{11}{13}\selectfont \textbf{Digital Onboarding — The Accountability Trap}}
		\begin{block}{\fontsize{7}{9}\selectfont \textbf{The Hard Question — CAMS Exam Favorite}}
			\scriptsize
			Which benefits are associated with digital onboarding technology according to FATF? (Select Two.)
		\end{block}
		\vspace{0.02cm}
		\answerbox{\large \textbf{1. Greater accountability in compliance processes. 2. Lower cost and increased auditability.}}
		\vspace{0.02cm}
		\begin{block}{\fontsize{7}{9}\selectfont \textbf{Natural Explanation — Why You Got This Wrong}}
			\scriptsize
			\begin{itemize}
				\item \wronganswer{Elimination of human review requirements for high-risk customers} — \textbf{NO.} High-risk still requires human review.
				\item \textbf{Accountability:} Digital records create audit trails.
				\item \textbf{Lower cost:} Automation reduces manual effort.
				\item \textbf{Increased auditability:} Digital records are easier to track.
				\item \textbf{Key principle:} Digital onboarding = \textbf{accountability + cost savings + auditability}.
				\item \textbf{Exam Trap:} "Elimination of human review" is incorrect.
			\end{itemize}
		\end{block}
		\vspace{0.02cm}
		\recallbox{\scriptsize \textbf{Key Rule:} Digital onboarding = \textbf{accountability + lower cost + auditability}.}
	\end{frame}
	
	% ================================================================
	% SLIDE 29: AML-SENSITIVE COMMUNICATIONS
	% ================================================================
	\begin{frame}
		\frametitle{\fontsize{11}{13}\selectfont \textbf{AML Communications — The Pre-Approved Scripts Trap}}
		\begin{block}{\fontsize{7}{9}\selectfont \textbf{The Hard Question — CAMS Exam Favorite}}
			\scriptsize
			A bank is developing new training materials for customer-facing staff regarding AML-sensitive communications. What should the materials emphasize to reduce risk? (Select Two.)
		\end{block}
		\vspace{0.02cm}
		\answerbox{\large \textbf{1. Avoid mentioning any internal reviews, even if prompted by the customer. 2. Use pre-approved scripts when addressing service disruptions caused by internal reviews.}}
		\vspace{0.02cm}
		\begin{block}{\fontsize{7}{9}\selectfont \textbf{Natural Explanation — Why You Got This Wrong}}
			\scriptsize
			\begin{itemize}
				\item \wronganswer{Refer customers to the AML department for financial advice} — \textbf{NO.} AML department doesn't give financial advice.
				\item \textbf{Avoid mentioning internal reviews:} Could tip off customers.
				\item \textbf{Pre-approved scripts:} Ensures consistency and reduces risk.
				\item \textbf{Key principle:} AML communications = \textbf{avoid tipping off + use scripts}.
				\item \textbf{Exam Trap:} "Refer to AML for advice" is incorrect — AML is for compliance, not advice.
			\end{itemize}
		\end{block}
		\vspace{0.02cm}
		\recallbox{\scriptsize \textbf{Key Rule:} AML communications = \textbf{avoid mentioning reviews + use pre-approved scripts}.}
	\end{frame}
	
	% ================================================================
	% SLIDE 30: ALERT CLOSURE STANDARDS
	% ================================================================
	\begin{frame}
		\frametitle{\fontsize{11}{13}\selectfont \textbf{Alert Closure — The Documentation Trap}}
		\begin{block}{\fontsize{7}{9}\selectfont \textbf{The Hard Question — CAMS Exam Favorite}}
			\scriptsize
			A compliance team aims to enhance its standards for determining when enough research has been done to close or escalate an alert. Which practices should be included? (Select Three.)
		\end{block}
		\vspace{0.02cm}
		\answerbox{\large \textbf{1. Documenting all steps taken and data sources reviewed. 2. Supplementing decisions with verified beneficial ownership documentation. 3. Establish an audit timeline that another independent reviewer could rely on.}}
		\vspace{0.02cm}
		\begin{block}{\fontsize{7}{9}\selectfont \textbf{Natural Explanation — Why You Got This Wrong}}
			\scriptsize
			\begin{itemize}
				\item \wronganswer{Closing alerts based on analyst discretion without documentation} — \textbf{NO.} Must be documented.
				\item \textbf{Document all steps:} Shows what was done.
				\item \textbf{Beneficial ownership:} Verified documentation supports decisions.
				\item \textbf{Audit timeline:} Another reviewer could replicate the process.
				\item \textbf{Key principle:} Alert closure = \textbf{documentation + verification + auditability}.
				\item \textbf{Exam Trap:} "Analyst discretion" without documentation is unacceptable.
			\end{itemize}
		\end{block}
		\vspace{0.02cm}
		\recallbox{\scriptsize \textbf{Key Rule:} Alert closure = \textbf{documentation + verification + auditability}.}
	\end{frame}
	
	% ================================================================
	% SLIDE 31: DEEPFAKE PREVENTION TOOLS
	% ================================================================
	\begin{frame}
		\frametitle{\fontsize{11}{13}\selectfont \textbf{Deepfake Prevention — The Texture Trap}}
		\begin{block}{\fontsize{7}{9}\selectfont \textbf{The Hard Question — CAMS Exam Favorite}}
			\scriptsize
			A fraudster used a high-resolution image and a deepfake video to bypass facial recognition verification. What tools could have prevented this fraud attempt? (Select Two.)
		\end{block}
		\vspace{0.02cm}
		\answerbox{\large \textbf{1. Use of software that can detect inconsistencies in light or skin texture. 2. Prompts for the customer to blink, smile, or repeat a phrase.}}
		\vspace{0.02cm}
		\begin{block}{\fontsize{7}{9}\selectfont \textbf{Natural Explanation — Why You Got This Wrong}}
			\scriptsize
			\begin{itemize}
				\item \wronganswer{Real-time fraud risk scoring algorithms using behavioral analytics} — \textbf{NO.} That's behavioral, not deepfake detection.
				\item \textbf{Texture analysis:} Detects inconsistencies in light, skin texture, shadows.
				\item \textbf{Active liveness:} Blink, smile, repeat phrase — requires real-time action.
				\item \textbf{Key principle:} Deepfake prevention = \textbf{texture analysis + active liveness}.
				\item \textbf{Exam Trap:} "Behavioral analytics" is different from deepfake detection.
			\end{itemize}
		\end{block}
		\vspace{0.02cm}
		\recallbox{\scriptsize \textbf{Key Rule:} Deepfake prevention = \textbf{texture analysis + active liveness prompts}.}
	\end{frame}
	
	% ================================================================
	% SLIDE 32: AFC DATA EXTRACTION CHALLENGE
	% ================================================================
	\begin{frame}
		\frametitle{\fontsize{11}{13}\selectfont \textbf{Data Extraction — The Normalization Trap}}
		\begin{block}{\fontsize{7}{9}\selectfont \textbf{The Hard Question — CAMS Exam Favorite}}
			\scriptsize
			Which best describes a challenge during AFC data extraction?
		\end{block}
		\vspace{0.02cm}
		\answerbox{\large \textbf{Data may come in different formats, requiring normalization during integration.}}
		\vspace{0.02cm}
		\begin{block}{\fontsize{7}{9}\selectfont \textbf{Natural Explanation — Why You Got This Wrong}}
			\scriptsize
			\begin{itemize}
				\item \wronganswer{Data is always in a consistent format} — \textbf{NO.} Data comes in many formats.
				\item \textbf{Different formats:} PDF, JSON, XML, CSV, etc.
				\item \textbf{Normalization:} Converting to a consistent format for processing.
				\item \textbf{Key principle:} Data extraction challenge = \textbf{different formats requiring normalization}.
				\item \textbf{Exam Trap:} "Consistent format" is unrealistic — data is varied.
			\end{itemize}
		\end{block}
		\vspace{0.02cm}
		\recallbox{\scriptsize \textbf{Key Rule:} Data extraction = \textbf{normalization of different formats}.}
	\end{frame}
	
	% ================================================================
	% SLIDE 33: DATA MINING REQUIREMENTS
	% ================================================================
	\begin{frame}
		\frametitle{\fontsize{11}{13}\selectfont \textbf{Data Mining — The Intelligent Filters Trap}}
		\begin{block}{\fontsize{7}{9}\selectfont \textbf{The Hard Question — CAMS Exam Favorite}}
			\scriptsize
			Which are requirements for effective data mining and matching in AFC controls? (Select Two.)
		\end{block}
		\vspace{0.02cm}
		\answerbox{\large \textbf{1. Comparing extracted data against structured and unstructured risk sources. 2. Applying intelligent filters to reduce false positives during entity resolution.}}
		\vspace{0.02cm}
		\begin{block}{\fontsize{7}{9}\selectfont \textbf{Natural Explanation — Why You Got This Wrong}}
			\scriptsize
			\begin{itemize}
				\item \wronganswer{Focusing on internal customer files for mining purposes} — \textbf{NO.} Need both internal and external sources.
				\item \textbf{Structured and unstructured sources:} Both are needed for effective matching.
				\item \textbf{Intelligent filters:} Reduce false positives in entity resolution.
				\item \textbf{Key principle:} Data mining = \textbf{both structured/unstructured + filters}.
				\item \textbf{Exam Trap:} "Internal files only" is incomplete — need external sources too.
			\end{itemize}
		\end{block}
		\vspace{0.02cm}
		\recallbox{\scriptsize \textbf{Key Rule:} Data mining = \textbf{structured + unstructured sources + intelligent filters}.}
	\end{frame}
	
	% ================================================================
	% SLIDE 34: RISK-BASED AUTHENTICATION
	% ================================================================
	\begin{frame}
		\frametitle{\fontsize{11}{13}\selectfont \textbf{Risk-Based Authentication — The Deviation Trap}}
		\begin{block}{\fontsize{7}{9}\selectfont \textbf{The Hard Question — CAMS Exam Favorite}}
			\scriptsize
			How do risk-based authentication systems optimize user experience?
		\end{block}
		\vspace{0.02cm}
		\answerbox{\large \textbf{They prompt additional checks only when user behavior deviates from expectations.}}
		\vspace{0.02cm}
		\begin{block}{\fontsize{7}{9}\selectfont \textbf{Natural Explanation — Why You Got This Wrong}}
			\scriptsize
			\begin{itemize}
				\item \wronganswer{They always require multi-factor authentication} — \textbf{NO.} That's not risk-based.
				\item \textbf{Risk-based:} Only steps up when needed.
				\item \textbf{Behavior deviation:} Additional checks only when behavior is unusual.
				\item \textbf{Key principle:} Risk-based authentication = \textbf{step up only when needed}.
				\item \textbf{Exam Trap:} "Always MFA" is not risk-based — that's standard authentication.
			\end{itemize}
		\end{block}
		\vspace{0.02cm}
		\recallbox{\scriptsize \textbf{Key Rule:} Risk-based authentication = \textbf{step up only when behavior deviates}.}
	\end{frame}
	
	% ================================================================
	% SLIDE 35: REGULATOR ASSESSMENT — FINTECH ONBOARDING
	% ================================================================
	\begin{frame}
		\frametitle{\fontsize{11}{13}\selectfont \textbf{Fintech Onboarding — The Verification Trap}}
		\begin{block}{\fontsize{7}{9}\selectfont \textbf{The Hard Question — CAMS Exam Favorite}}
			\scriptsize
			A regulator is assessing a fintech startup's customer identity verification program. The firm currently relies solely on device recognition technology. What additional measures could the fintech implement at initial customer onboarding? (Select Two.)
		\end{block}
		\vspace{0.02cm}
		\answerbox{\large \textbf{1. Implementing document verification with biometric authentication. 2. Integrating with national electronic identity verification systems.}}
		\vspace{0.02cm}
		\begin{block}{\fontsize{7}{9}\selectfont \textbf{Natural Explanation — Why You Got This Wrong}}
			\scriptsize
			\begin{itemize}
				\item \wronganswer{Enhancing digital footprint analysis through IP address verification} — \textbf{NO.} That's similar to device recognition, not additional.
				\item \textbf{Document verification:} ID document checks.
				\item \textbf{Biometric authentication:} Fingerprint, facial recognition.
				\item \textbf{National eID systems:} Government-backed identity verification.
				\item \textbf{Key principle:} Onboarding = \textbf{document verification + biometrics + national eID}.
				\item \textbf{Exam Trap:} "IP verification" is similar to device recognition — not a new measure.
			\end{itemize}
		\end{block}
		\vspace{0.02cm}
		\recallbox{\scriptsize \textbf{Key Rule:} Onboarding = \textbf{document verification + biometrics + national eID}.}
	\end{frame}
	
	% ================================================================
	% SLIDE 36: FACIAL RECOGNITION RISK
	% ================================================================
	\begin{frame}
		\frametitle{\fontsize{11}{13}\selectfont \textbf{Facial Recognition — The Irreplaceable Trap}}
		\begin{block}{\fontsize{7}{9}\selectfont \textbf{The Hard Question — CAMS Exam Favorite}}
			\scriptsize
			Which risk is specifically associated with facial recognition technology compared to traditional password-based authentication methods?
		\end{block}
		\vspace{0.02cm}
		\answerbox{\large \textbf{Compromised biometric data is difficult to change or reset unlike passwords.}}
		\vspace{0.02cm}
		\begin{block}{\fontsize{7}{9}\selectfont \textbf{Natural Explanation — Why You Got This Wrong}}
			\scriptsize
			\begin{itemize}
				\item \wronganswer{The technology requires more storage capacity than traditional methods} — \textbf{NO.} That's not a unique risk.
				\item \textbf{Passwords can be changed:} If compromised, reset them.
				\item \textbf{Biometrics can't be changed:} Can't get a new face.
				\item \textbf{Permanent compromise:} Once biometric data is stolen, it's forever.
				\item \textbf{Key principle:} Biometric risk = \textbf{irreplaceable}.
				\item \textbf{Exam Trap:} "Storage capacity" is not the unique risk.
			\end{itemize}
		\end{block}
		\vspace{0.02cm}
		\recallbox{\scriptsize \textbf{Key Rule:} Biometric risk = \textbf{can't change like passwords}.}
	\end{frame}
	
	% ================================================================
	% SLIDE 37: AUTHENTICATION TECHNOLOGIES
	% ================================================================
	\begin{frame}
		\frametitle{\fontsize{11}{13}\selectfont \textbf{Authentication — The Biometrics Trap}}
		\begin{block}{\fontsize{7}{9}\selectfont \textbf{The Hard Question — CAMS Exam Favorite}}
			\scriptsize
			Which technologies enhance authentication and security in financial systems? (Select Three.)
		\end{block}
		\vspace{0.02cm}
		\answerbox{\large \textbf{1. Multi-factor authentication. 2. Physiological biometrics. 3. Behavioral biometrics.}}
		\vspace{0.02cm}
		\begin{block}{\fontsize{7}{9}\selectfont \textbf{Natural Explanation — Why You Got This Wrong}}
			\scriptsize
			\begin{itemize}
				\item \wronganswer{Password biometrics} — \textbf{NO.} This doesn't exist.
				\item \textbf{MFA:} Multiple authentication factors.
				\item \textbf{Physiological biometrics:} Physical traits (fingerprint, face, iris).
				\item \textbf{Behavioral biometrics:} Behavior patterns (typing, mouse movements).
				\item \textbf{Key principle:} Biometrics = \textbf{physical or behavioral}. Passwords are knowledge-based.
				\item \textbf{Exam Trap:} "Password biometrics" is a trap — doesn't exist.
			\end{itemize}
		\end{block}
		\vspace{0.02cm}
		\recallbox{\scriptsize \textbf{Key Rule:} Authentication = \textbf{MFA + physiological + behavioral}.}
	\end{frame}
	
	% ================================================================
	% SLIDE 38: CONCENTRATION ACCOUNT RISK
	% ================================================================
	\begin{frame}
		\frametitle{\fontsize{11}{13}\selectfont \textbf{Concentration Accounts — The Traceability Trap}}
		\begin{block}{\fontsize{7}{9}\selectfont \textbf{The Hard Question — CAMS Exam Favorite}}
			\scriptsize
			During a routine compliance review, a financial institution finds that its treasury unit uses a concentration account that lacks sufficient oversight. What is the most likely regulatory concern?
		\end{block}
		\vspace{0.02cm}
		\answerbox{\large \textbf{The institution might inadvertently facilitate money laundering due to poor traceability.}}
		\vspace{0.02cm}
		\begin{block}{\fontsize{7}{9}\selectfont \textbf{Natural Explanation — Why You Got This Wrong}}
			\scriptsize
			\begin{itemize}
				\item \wronganswer{The account might be used for legitimate treasury operations} — \textbf{NO.} That's not a concern.
				\item \textbf{Concentration accounts:} Pool funds from multiple sources.
				\item \textbf{Poor traceability:} Can't track funds back to origin.
				\item \textbf{Money laundering risk:} Funds can be mixed and moved without trace.
				\item \textbf{Key principle:} Concentration accounts = \textbf{traceability risk}.
				\item \textbf{Exam Trap:} "Legitimate operations" ignores the risk.
			\end{itemize}
		\end{block}
		\vspace{0.02cm}
		\recallbox{\scriptsize \textbf{Key Rule:} Concentration accounts = \textbf{poor traceability = money laundering risk}.}
	\end{frame}
	
	% ================================================================
	% SLIDE 39: CLUSTERING BENEFIT
	% ================================================================
	\begin{frame}
		\frametitle{\fontsize{11}{13}\selectfont \textbf{Clustering Benefit — The Pattern Trap}}
		\begin{block}{\fontsize{7}{9}\selectfont \textbf{The Hard Question — CAMS Exam Favorite}}
			\scriptsize
			What is the primary benefit of clustering within transaction monitoring systems?
		\end{block}
		\vspace{0.02cm}
		\answerbox{\large \textbf{It reveals shared patterns among entities to detect potential financial crime.}}
		\vspace{0.02cm}
		\begin{block}{\fontsize{7}{9}\selectfont \textbf{Natural Explanation — Why You Got This Wrong}}
			\scriptsize
			\begin{itemize}
				\item \wronganswer{It eliminates all false positives} — \textbf{NO.} Clustering doesn't eliminate false positives.
				\item \textbf{Shared patterns:} Groups entities with similar behaviors.
				\item \textbf{Detects financial crime:} Identifies networks of suspicious activity.
				\item \textbf{Key principle:} Clustering = \textbf{shared patterns to detect crime}.
				\item \textbf{Exam Trap:} "Eliminates false positives" is unrealistic.
			\end{itemize}
		\end{block}
		\vspace{0.02cm}
		\recallbox{\scriptsize \textbf{Key Rule:} Clustering = \textbf{shared patterns for crime detection}.}
	\end{frame}
	
	% ================================================================
	% SLIDE 40: DATA MINING FUNCTIONS
	% ================================================================
	\begin{frame}
		\frametitle{\fontsize{11}{13}\selectfont \textbf{Data Mining Functions — The Comparison Trap}}
		\begin{block}{\fontsize{7}{9}\selectfont \textbf{The Hard Question — CAMS Exam Favorite}}
			\scriptsize
			What are key functions performed during AFC data mining and matching? (Select Three.)
		\end{block}
		\vspace{0.02cm}
		\answerbox{\large \textbf{1. Uncovering patterns in large complex datasets. 2. Comparing data from different sources to identify inconsistencies. 3. Detecting anomalies indicative of fraudulent behavior.}}
		\vspace{0.02cm}
		\begin{block}{\fontsize{7}{9}\selectfont \textbf{Natural Explanation — Why You Got This Wrong}}
			\scriptsize
			\begin{itemize}
				\item \wronganswer{Maintaining databases of known criminal activities} — \textbf{NO.} That's database maintenance, not mining.
				\item \textbf{Patterns:} Uncovering hidden relationships.
				\item \textbf{Comparison:} Identifying inconsistencies across sources.
				\item \textbf{Anomalies:} Detecting unusual behavior.
				\item \textbf{Key principle:} Data mining = \textbf{patterns + comparison + anomalies}.
				\item \textbf{Exam Trap:} "Maintaining databases" is operational, not mining.
			\end{itemize}
		\end{block}
		\vspace{0.02cm}
		\recallbox{\scriptsize \textbf{Key Rule:} Data mining = \textbf{patterns + comparison + anomalies}.}
	\end{frame}
	
	% ================================================================
	% SLIDE 41: FALSE NEGATIVES — DATA QUALITY
	% ================================================================
	\begin{frame}
		\frametitle{\fontsize{11}{13}\selectfont \textbf{False Negatives — The Data Quality Trap}}
		\begin{block}{\fontsize{7}{9}\selectfont \textbf{The Hard Question — CAMS Exam Favorite}}
			\scriptsize
			An internal review finds a high rate of false negatives in the organization's monitoring system due to outdated KYC files and inconsistent country codes. Which is the most reasonable conclusion?
		\end{block}
		\vspace{0.02cm}
		\answerbox{\large \textbf{Data quality controls are insufficient and must be reinforced at the source level.}}
		\vspace{0.02cm}
		\begin{block}{\fontsize{7}{9}\selectfont \textbf{Natural Explanation — Why You Got This Wrong}}
			\scriptsize
			\begin{itemize}
				\item \wronganswer{The monitoring system should be replaced} — \textbf{NO.} The issue is data quality, not the system.
				\item \textbf{Outdated KYC:} Source data is not reliable.
				\item \textbf{Inconsistent country codes:} Data is not standardized.
				\item \textbf{Key principle:} False negatives = \textbf{data quality issue at source}.
				\item \textbf{Exam Trap:} "Replace system" ignores the root cause — data quality.
			\end{itemize}
		\end{block}
		\vspace{0.02cm}
		\recallbox{\scriptsize \textbf{Key Rule:} False negatives = \textbf{data quality controls insufficient at source}.}
	\end{frame}
	
	% ================================================================
	% SLIDE 42: FINAL EXAM TIPS — DOMAIN D
	% ================================================================
	\begin{frame}
		\frametitle{\fontsize{11}{13}\selectfont \textbf{Domain D — Complete Exam Traps Summary}}
		\scriptsize
		\begin{block}{\fontsize{7}{9}\selectfont \textbf{All Hard Questions — Natural Explanations}}
			\scriptsize
			\begin{enumerate}
				\item \textbf{Sanctions Lists:} Update because individuals change aliases/affiliations.
				\item \textbf{Transaction Monitoring:} Improve with behavioral + peer-group comparisons.
				\item \textbf{Dynamic Data:} Real-time geolocation + live adverse media feeds.
				\item \textbf{Governance:} Align with regulatory expectations.
				\item \textbf{External Data:} Continuous pipeline + standardization.
				\item \textbf{Risk Scoring:} New flags = possible success at identifying missed risk.
				\item \textbf{KYC Onboarding:} Automated data entry + document verification.
				\item \textbf{OSINT/AI Documentation:} Feature selection + training data + output explanations.
				\item \textbf{Clustering:} Shared behavior + cross-transaction connections.
				\item \textbf{Human-in-the-Loop:} Expert assessment + explainability.
				\item \textbf{Entity Resolution:} Reveals hidden connections.
				\item \textbf{Behavioral Monitoring Governance:} Thresholds vs. SARs + tuning reviews.
				\item \textbf{OSINT Quality:} Reassess sources + document reliability.
				\item \textbf{Risk-Based Tech Selection:} Fit for risk + vendor eval + strategy alignment.
				\item \textbf{ML Benefits:} Continuous learning + improved accuracy.
				\item \textbf{Card Risks:} Anonymity + rapid transactions + high-risk customers.
				\item \textbf{Data Gaps:} Missing identifiers + inconsistent formats.
				\item \textbf{DNFBP Supervision:} Extend AML/CFT to lawyers, accountants, real estate.
				\item \textbf{PSP Strategy:} Customize per jurisdiction and product.
				\item \textbf{Prepaid Product:} Reassess risks + adjust thresholds FIRST.
				\item \textbf{Data Visualization:} Faster resolution + more false negatives = overlooking signals.
				\item \textbf{Data Validation:} Accuracy + functionality + appropriate alerts.
				\item \textbf{Data Lineage:} Transformations history + source of each field.
				\item \textbf{Real-Time Detection:} ML trained on known suspicious behavior.
				\item \textbf{Data Standardization:} Data dictionary + mapping schema.
				\item \textbf{Digital Onboarding:} Accountability + lower cost + auditability.
				\item \textbf{AML Communications:} Avoid mentioning reviews + pre-approved scripts.
				\item \textbf{Alert Closure:} Documentation + verification + auditability.
				\item \textbf{Deepfake Prevention:} Texture analysis + active liveness prompts.
				\item \textbf{Data Extraction:} Normalization of different formats.
				\item \textbf{Data Mining:} Structured + unstructured sources + intelligent filters.
				\item \textbf{Risk-Based Authentication:} Step up only when behavior deviates.
				\item \textbf{Fintech Onboarding:} Document verification + biometrics + national eID.
				\item \textbf{Facial Recognition Risk:} Can't change like passwords.
				\item \textbf{Authentication:} MFA + physiological + behavioral.
				\item \textbf{Concentration Accounts:} Traceability risk.
				\item \textbf{Clustering Benefit:} Shared patterns for crime detection.
				\item \textbf{Data Mining Functions:} Patterns + comparison + anomalies.
				\item \textbf{False Negatives:} Data quality controls insufficient at source.
			\end{enumerate}
		\end{block}
		\vspace{0.02cm}
		\recallbox{\scriptsize \textbf{Final Rule:} Domain D tests \textbf{technology application} in AML. Understand \textbf{what each tool does} and \textbf{how it helps}.}
	\end{frame}
	
	% ================================================================
	% END SLIDE
	% ================================================================
	\begin{frame}
		\centering
		\vspace{1.5cm}
		{\fontsize{18}{22}\selectfont \textbf{Domain D — Complete}}\\[0.2cm]
		{\fontsize{11}{13}\selectfont Tools \& Technologies to Fight Financial Crime}\\[0.1cm]
		\rule{4cm}{0.5pt}\\[0.1cm]
		{\fontsize{8}{10}\selectfont AI/ML Tools · RPA · Digital Onboarding · Data Validation · Data Lineage}\\[0.05cm]
		{\fontsize{8}{10}\selectfont Blockchain Analytics · Entity Resolution · Clustering · Batch Screening}\\[0.05cm]
		{\fontsize{8}{10}\selectfont Device Intelligence · Biometrics · Liveness Detection · Perpetual KYC}\\[0.05cm]
		{\fontsize{8}{10}\selectfont OSINT · Network Analysis · Cryptocurrency · Deepfake Detection}\\[0.1cm]
		\rule{4cm}{0.5pt}\\[0.1cm]
		{\fontsize{8}{10}\selectfont \textit{All traps explained. All concepts covered. Ready for the exam.}}
		\vspace{0.5cm}
	\end{frame}
	
\end{document}